本文提出了一种面向资源受限电池管理系统的轻量化锂离子电池SOH估计框架，以缓解健康指标跨电池稳定性不足以及预测精度与计算效率难以兼顾的问题。研究从系统性健康指标构建和轻量化网络设计两个方面展开。所提出的MS-CCCT基于多节特征开发电池的PCC与SCC，对恒流充电电压窗口进行组级多尺度自适应标定；在此基础上，采用PCC/SCC双阈值准入与双相关冗余约束确定模型输入。在保持特征定义与参数不变的条件下，最终入选指标在相应独立评价电池上仍与SOH保持较强的线性和单调关联。在序列建模方面，MS-AgentNet将ReLU²智能体注意力与小核和大核深度可分离卷积相结合，以协同表征局部退化变化、跨位置全局信息和长尺度退化趋势。在智能体数量固定时，RAA利用少量可学习静态智能体完成信息聚合与广播，将注意力相关性交互的理论复杂度由$O(N^2d)$降低至$O(Nn_a d)$。

在Oxford、CALCE CS2、CALCE CX2和MIT/Severson四个公开电池数据集上的实验结果表明，与CNN-Transformer、CNN-LSTM、Transformer和LSTM四种基线模型相比，MS-AgentNet在四个数据集的独立评价结果中均取得了最佳综合表现。具体而言，在CALCE CX2数据集上，其平均绝对百分比误差相比CNN-Transformer降低了59.75\%。此外，MS-AgentNet具有较低的参数量和存储开销，体现了其轻量化优势。消融实验与卷积尺度补充实验显示，多尺度深度可分离卷积与RAA的互补融合能够提升模型在不同退化场景下的综合表现。总体来看，所提框架在SOH估计精度与资源开销之间取得了良好平衡，并在各数据集的独立评价电池上表现出较好的域内跨电池泛化能力，具备应用于资源受限电池管理系统的潜力。

尽管取得了上述结果，所提框架仍依赖可辨识的充放电数据片段，且其跨数据集适应能力会受到数据域差异、源域选择和目标域数据量的影响。未来将重点研究面向不完整充电和动态运行片段的健康指标构建方法及轻量化域适应策略，并在更多电池体系、实际车辆运行数据和嵌入式硬件平台上开展验证，以进一步提升该框架的工程适用性与部署可靠性。
